

Through this report, we have looked at the discourse around the eucalyptus tree. We have seen that there is a spike in the number of documents related to the tree around 1870, driven mostly by the generalist press, which cannot be fully explained by the general rise of published documents, thus showing a specific phenomenon happening at the time. We also discussed that these documents can be analyzed through a topic model, which, together with the plotting of documents per place of publication, reveals interesting key studies, such as Lyon or Madagascar. Moreover, it allows the analysis of specific themes, such as the topic of malaria. The latter which reveals that the topic of eucalyptus not only appears in the contexts where it might be expected (generalist press, monographs or official documents), but also through religious journals, which relay the work of eucalyptus plantations around Rome.

These results, while confirming intuitions about the discourse, still show some light on interesting ways it was shaped. Paris, being the cultural center of the francophone world is the city that really drives the subject. It is basically the only place that does publish about the plant before 1870 in the corpus, while still being largely ahead of the other regions after that date. Moreover, it is the center of any scientific production, publishing no less than 80\% of all monographs of the corpus. However, we cannot neglect the contribution of the documents published in the regions of plantation, which also help shape the discourse locally. This is particularly striking in the topic distribution between the south of France and the colonies. By comparing the two case studies of Lyon and Madagascar, we see a clear difference in the topics mentioned as well as in the actors at play. While in the mainland it manifests mainly through advertisements, in Madagascar it is through the colonial institutions that eucalyptus appear in the documents (and the same phenomenon seems to occur in North Africa). Indeed, we see that the properties linked with the plant seem to not be used (or at least not advertised) in the colonies, while being heavily done so in the mainland. This can be seen through the disproportion between medicinal meta-topics (diseases, medicine and medication and chemical properties), which represent less than 5\% of documents in the colonial regions, while often more than 15\% in the mainland. We can thus easily see the
domination of the mainland over its colonial empire. Paris and Lyon do use the plant not only for its drainage properties but also for its medicinal ones, while Madagascar and North Africa are only subjected to the former. 

This conclusion can be linked with prior research on the tree. Eucalyptus as a colonial tool to "correct" the environment, so that it is better suited to the colonial interests, has been discussed before by Davin-Mortier\footnote{\cite[3, 6]{davin-mortier_leucalyptus_2026}}, showing that it was used to sanitize marshes that could then be used, for example, for food cultivation or wood production. Such policies, marked by the high amount of institutional topics, do seem to be present in the regions presented here (North Africa, Madagascar). These studies being qualitative in nature, this project allows instead for a wider quantitative analysis of the subject, which does demonstrate a difference in the discourse between the colonizer and colonized regions. 

Finally, we also have to discuss the numerous limitations of the study. The biggest one, as mentioned in Section \ref{sec:met_files} is related to the filtering of the files. Since Gallica's API limitations did not allow the download of all potential files, a filtering had to be put in place. However, the way it was implemented, while sounding relevant at the time, produced some biases related to which journals were or were not kept, as well as removing a large portion of the generalist press. The cracks show for example when examining the biggest publishing cities in the south of France, where places such as "Toulouse", "Marseille" or "Bordeaux" do not show up in the top 5, whereas smaller towns do so. However, the regions results tend to show that a high amount of the discourse is used through advertisements, and it is probably the case for these cities as well. Nonetheless, a further study should take these places into account, so that a less biased corpus can be formed. Speaking of biases, we also need to take into account the biases deriving from the original Gallica database. While more abstract, the digitization process does involve some biases in which documents are digitized, how their metadata fields are filled out and more importantly which documents from the colonized regions are included. This might play a role in the amount of institutional topics in these regions, as the French libraries contributing to the database might have the more official and colonial documents, while not necessarily having access to or not having digitized the more local newspapers. This can also be linked with the language of the study, as these local newspapers would not all have been published in French, thus not entering this corpus. Then, as with any similar project, the OCR presents some limitations. The model used, not necessarily being suited for certain documents types such as the general press, makes it so that properly isolating the relevant context is a hard task to generalize. To mitigate this, we have chosen to take into account a very small context in order to take into account this issue. For a further analysis however, it would be recommended to not download Gallica's own text, but the page's scan and perform the OCR in-house, hopefully capturing better the structure of such documents. Finally, while performing a quantitative analysis of the discourse is useful to understand at a macro scale the dynamics and tendencies of a large corpus, it does not detract for the need of an in-depth qualitative analysis of the documents. While we tried in places to provide context directly from the documents, further studies should go more in-depth in this disrection. Notably, this analysis can provide the basis for discovering such documents that might have not been studied before. 